%%
%
% ARQUIVO: dados-seminario.tex
%
% VERSÃO: 1.1
% DATA: Abril de 2016
% AUTOR: Coordenação PPgSC
% 
%  Arquivo tex com os dados acerca do relatório de Avaliação de Trabalho de Dissertação.
%
%
%%

%%% AUTOR DA APRESENTAÇÃO DE SEMINÁRIO (Nome completo)
\autor{Nycolas Wenderson da Silva dos Santos}

%%% CÓDIGO DO AUTOR DA APRESENTAÇÃO DE SEMINÁRIO
\codigoautor{SC 23202}

%%% POSTO DO AUTOR DA APRESENTAÇÃO DE SEMINÁRIO
% ---
%  se o autor é CIVIL, REMOVA ESTA LINHA
%  se o autor é MILITAR, coloque CORRETAMENTE o POSTO aqui
% ---
%\postoautor{1 Ten}

%%% TITULO DO TRABALHO DE DISSERTAÇÃO
\titulo{ATRIBUIÇÃO DE ALGORITMOS DE CRIPTOGRAFIA LEVE EM CENÁRIO
CIPHERTEXT-ONLY VIA APRENDIZADO DE MÁQUIN}

%%% TÍTULO ANTERIOR DO TRABALHO DE DISSERTAÇÃO
% ---
%  se o título do trabalho NÃO foi alterado desde a última
%     apresentação (Proposta ou Seminário), REMOVA ESTA LINHA
%  se o título do trabalho FOI ALTERADO, coloque aqui o ANTIGO
%     título COMPLETO do trabalho:
% ---
%\titantigo{Título Anterior Completo do Trabalho de Dissertação}

%%% DATA DA APRESENTAÇÃO DE SEMINÁRIO (formato {dd}{Mmmmm}{aaaa})
\dataapresentacao{07}{Julho}{2026}

%%% ÁREA DE CONCENTRAÇÃO DO TRABALHO DE DISSERTAÇÃO
% ---
%  VER SITE do PPGSC para ver as Áreas de Concentração do Programa.
%  Em 2016 só existe uma Área de Concentração:
%     Ciência da Computação
% ---
\area{Ciência da Computação}

%%% LINHA DE PESQUISA DO TRABALHO DE DISSERTAÇÃO
% ---
%  VER SITE do PPGSC para ver as Linhas de Pesquisa do Programa.
%  Em 2016 só existem três Linhas de Pesquisa:
%     Metodologia da Computação
%     Sistemas de Computação
%     Engenharia de Sistemas e Informação
% ---
\linha{Sistemas de Computação}

%%% ORIENTADOR DO TRABALHO DE DISSERTAÇÃO
% ---
%  CAMPO 1: Nome completo
%  CAMPO 2: D (para D.Sc.); P (para Ph.D.); ou qualquer coisa (inclusive VAZIO) - o que for escrito aparecerá no documento
% ---
\orientador{Jose Antonio Moreira Xéxeo}{D}

%%% POSTO DO ORIENTADOR DO TRABALHO DE DISSERTAÇÃO
% ---
%  se o orientador é CIVIL, REMOVA ESTA LINHA
%  se o orientador é MILITAR, coloque CORRETAMENTE o POSTO aqui:
%     pode ser: 1 Ten; Cap; Maj; Ten Cel; Cel
% ---
\postoorientador{Ten Cel}

%%% CO-ORIENTADOR DO TRABALHO DE DISSERTAÇÃO
% ---
%  se NÃO HOUVER co-orientador, REMOVA ESTA LINHA
%  preenchimento idêntico a \orientador{}{}
% ---
%\coorientador{Nome Completo do Co-orientador}{P}

%%% POSTO DO CO-ORIENTADOR DO TRABALHO DE DISSERTAÇÃO
% ---
%  se NÃO HOUVER co-orientador, REMOVA ESTA LINHA
%  caso contrário
%    se o co-orientador é CIVIL, REMOVA ESTA LINHA
%    se o co-orientador é MILITAR, coloque CORRETAMENTE o POSTO aqui:
%       pode ser: 1 Ten; Cap; Maj; Ten Cel; Cel
% ---
%\postocoorientador{Cap}

%%% COORDENADOR DE PÓS-GRADUAÇÃO
% ---
%  CAMPO 1: Nome completo
%  CAMPO 2: D (para D.Sc.); P (para Ph.D.); ou qualquer coisa (inclusive VAZIO) - o que for escrito aparecerá no documento
% ---
\coord{Gabriela Moutinho de Souza Dias}{P}

%%% POSTO DO COORDENADOR DE PÓS-GRADUAÇÃO
% ---
%  se o Coordenador de PG é CIVIL, REMOVA ESTA LINHA
%  se o Coordenador de PG é MILITAR, coloque CORRETAMENTE o POSTO aqui:
%     pode ser: 1 Ten; Cap; Maj; Ten Cel; Cel
% ---
\postocoord{Maj}

%%% CHEFE DA SEÇÃO (SE/8)
\chefe{Julio Cesar Duarte}

%%% POSTO DO CHEFE DA SEÇÃO (SE/8)
% ---
%  se o Chefe é CIVIL, REMOVA ESTA LINHA
%  se o Chefe é MILITAR, coloque CORRETAMENTE o POSTO aqui:
%     pode ser: 1 Ten; Cap; Maj; Ten Cel; Cel
% ---
\postochefe{Cel}
